\chapter*{摘要}

延迟（delay）是大多数动力系统固有的特性。除在时间上推移过程外，延迟还会显著影响系统性能。因此，研究延迟并对其加以考量通常很有价值。序贯决策问题（如\glspl{mdp}）也是动力系统，同样会受到延迟的影响，这并不令人意外。这些过程是\gls{rl}的基础框架，该范式的目标是创建能够通过与环境交互学习最大化效用的智能体（agent）。

\gls{rl}已取得强大、有时甚至惊人的实证结果，但延迟很少被明确纳入考量。关于延迟对\gls{mdp}影响的理解仍很有限。本论文旨在研究智能体在状态观测（state observation）或动作执行（action execution）中的延迟。不断变换视角以揭示问题的结构和特性。考虑多种延迟，并提出潜在解决方案。本文还旨在建立\gls{rl}文献中的著名框架与延迟框架之间的联系。因此，聚焦以下四个视角。

首先，考虑常值延迟（constant delay）。采用心理学启发的方法，研究预测近期未来以估计智能体动作影响的作用。强调该方法与\gls{rl}文献中模型的关系。借此机会，正式证明一个看似显而易见的事实：延迟越长，性能越低。实验分析最终验证该方法的有效性。

作为第二视角，考虑在延迟环境中模仿无延迟专家行为的简单方法。延迟最初保持恒定，但之后扩展研究到更特殊的延迟类型，如随机延迟（stochastic delay）。尽管方法简单，但展示了其强大的理论保证和实证结果。

第三视角同样针对常值延迟，考虑采用非平稳（non-stationary）无记忆行为。虽然看似忽略延迟，但该方法将延迟影响视为引导非平稳性的未观测变量。基于这一思想，提出一种具有理论基础的算法来学习此类行为，并在实际场景中测试。

最后，第四个视角考虑比常值延迟更广泛的模型，常值延迟作为其特例。该模型允许动作影响环境的多个未来转移（transition）。研究其理论特性以理解其特殊性。基于这些特性，一些\gls{rl}算法被排除，另一些则在多种实证研究中得到检验。作为更通用的延迟模型，对其理解对前几章中的常值延迟框架具有启示意义。
